\scanpage{16}
\mspage{\handwritten{28}}
% Oblique line from head of sheet to the bottom correction.
\deleted{\uline{Lemme} 5.1\deleted{\unsure{1}}\handwritten{2}. Soient $S$ un préschéma, $R$ un préschéma constant t\deleted{\unsure{r}}o\deleted{\unsure{o}}rdu sur $S$. Alors tout sous-préschéma fermé \deleted{\illegible{1}} \inserted{\handwritten{$Z$}} de $R$ qui est \deleted{de type fini}\inserted{quasi-compact} sur $S$, est fini sur $S$.

En effet, on est ramené au cas où $R$ est constant, donc de la forme $I_S$, où $I$ est un ensemble, donc réunion filtrante croissante des $J_S$, où $J$ parcourt les parties finies de $I$. On peut supposer de plus $S$ affine, alors \deleted{\unsure{$R_i$}} $Z$ est quasi-compact, donc contenu dans un des $J_S$. Comme $J_S$ est fini sur $S$, il en est de même de $Z$.

Le lemme 5.1\deleted{\unsure{1}}\handwritten{2}. établit déjà l'assertion entre parenthèses dans 5.\deleted{10}\handwritten{11}. (ii). \deleted{L'équivalence (ii) \quad (iii) est} \inserted{L'implication (ii)$\handwritten{\Longrightarrow}$(iii), et l\unsure{1|l}implication inverse} dans le cas $S$ localement noethérien est alors claire. Prouvons que (i)$\handwritten{\Longrightarrow}$(ii), pour ceci, soit $S'\handwritten{\longrightarrow}S$ un morphisme étale fini surjectif qui splitte $H$ donc $R$, de sorte que $R'$ s'écrit \deleted{comme} $M_S=\deleted{\unsure{R_m}}\handwritten{\coprod}_{m\handwritten{\in}M}R'_m$, où \deleted{\unsure{sont|tous|tout}} les $R'_m$ sont des ouverts disjoints de $R'$, \inserted{$S'$-}isomorphes à $S'$. Soit $g:R'\handwritten{\to}R$ la projection, qui est un morphisme fini étale surjectif, \inserted{donc un morphisme ouvert et fermé\unsure{;};} alors les $R_m=g(R'_m)$ sont des parties à \deleted{a}la fois ouvertes et fermées de $R$, et évidemment quasi-compacts sur $S$ puisque les $R_m$ le sont. \deleted{Enfin, prouvons (ii) \quad (i)}

Enfin, prouvons (ii)$\handwritten{\Longrightarrow}$(i) lorsque $H$ est de type fini sur $S$\handwritten{,} \inserted{\handwritten{$S$ quasi compact ou connexe.}} Soit $M$ son type, on peut écrire $M=\deleted{\unsure{T_{\oplus}|Z_{\oplus}}}\underline Z^r\times N$, où $r$ est un entier $\handwritten{\geq}0$ et $N$ un groupe fini. Considérons le préschéma
\[
P=\underline{\operatorname{Isom}}_{S\text{-gr}}(H,G)\handwritten{\subset}\underline{\operatorname{Hom}}_{S\text{-gr}}(H,G)\handwritten{\unsure{=|\Rightarrow}}Q
\]
où $G=D_S(M)$, cf 5.\handwritten{8}. et 5.\handwritten{10}. \deleted{Alors $Q$ est isomorphe} On a un isomorphis
\[
Q\handwritten{\simeq}\underline{\operatorname{Hom}}_{S\text{-gr}}(G,\underline G_m^r)\times\underline{\operatorname{Hom}}_{S\text{-gr}}(G,N)\handwritten{\simeq}R^r\times E\quad,
\]
où $E=\underline{\operatorname{Hom}}_{S\text{-gr}}(G,N)$ est fini sur $S$. Il s'ensuit que $Q$ est \deleted{somme de}\inserted{réunion de sous-}préschémas à la fois ouverts et fermés $Q_i$ finis sur $S$. \deleted{De même, Il s'ensuit que} \inserted{\handwritten{Donc}} $P$ est réunion \deleted{\unsure{d'un|disj}} de \deleted{\unsure{présc|prés}} sous-préschémas à la fois ouverts et fermés $P_i$\inserted{\handwritten{$=P\cap Q_i$}}} finis sur $S$. \deleted{L}Comme ils sont étales sur $S$, leurs
% archive sheet number: 15

